% ==============================================================================
% MÓDULO CAPÍTULO II: ARTICULACIÓN DIALÉCTICA CON LOS CUATRO EXISTENCIALES
% Archivo: cap_2/2.6_articulacion_lebenswelt.tex
% Enfoque: 100% Calidad de Atención (Donabedian y Watson)
% ==============================================================================
\subsection{Articulación dialéctica del escenario físico-social con las dimensiones existenciales del usuario (\textit{Lebenswelt})}
\label{sec:2.6_articulacion_existenciales}

En concordancia con los postulados epistemológicos de Max van Manen \cite{vanmanen1990}, John Creswell \cite{creswell2013} y Roberto Hernández-Sampieri \cite{sampieri2018}, articulados con el modelo del Proceso Interpersonal de Avedis Donabedian \cite{donabedian2005} y la Teoría del Cuidado Humano de Jean Watson \cite{watson2008}, el mundo de la vida (\textit{Lebenswelt}) de los usuarios se manifiesta a través de cuatro existenciales universales donde se construye el significado de la calidad de atención:

\begin{table}[H]
\centering
\small
\caption{Articulación dialéctica entre el contexto empírico de Belén y los existenciales de la calidad de atención (\textit{Lebenswelt}).}
\label{tab:articulacion_existenciales}
\vspace{0.2cm}
\begin{tabularx}{\textwidth}{p{3.2cm}p{4.8cm}X}
\toprule
\textbf{Dimensión Existencial (\textit{Lebenswelt})} & \textbf{Manifestación Empírica en el Contexto de Belén} & \textbf{Significado Vivencial de la Calidad Construido por el Usuario} \\
\midrule
\textbf{Relacionalidad} \newline (\textit{Relación vivida con los otros}) & 
Encuentro intersubjetivo entre el usuario quechua-mestizo vulnerable y el personal sanitario; presencia del saludo cordial (\textit{napaykuy}) o su omisión; mirada a los ojos frente a la mirada fija en el monitor; paciencia explicativa o respuestas cortantes. & 
La calidad se define en la dignidad del trato interpersonal: el usuario busca acogida humana (\textit{allin chaskiy}), respeto mutuo (\textit{respetanakuy}) y comprensión en su propia lengua originaria. La calidez del cuidado de enfermería transforma la vulnerabilidad en alivio, confianza y gratitud. \\
\addlinespace
\textbf{Temporalidad} \newline (\textit{Tiempo vivido}) & 
Puntualidad en el inicio de la atención; duración efectiva del contacto clínico en consultorio; lapsos de espera frente a una consulta médica de apenas 10 a 15 minutos; ritmo pausado para escuchar dudas. & 
Vivencia de la desvalorización del tiempo del paciente frente a la prisa protocolar. Una atención de calidad se experimenta cuando el profesional no muestra apuro, escucha sin interrumpir y dedica tiempo suficiente para clarificar el tratamiento. \\
\addlinespace
\textbf{Espacialidad} \newline (\textit{Espacio vivido}) & 
Consultorios con tabiquería delgada que deja filtrar conversaciones clínicas; pasadizos congestionados; puertas entreabiertas durante la exploración médica; comodidad y limpieza del consultorio. & 
Sensación de exposición pública o resguardo de la intimidad. La calidad espacial exige consultorios cerrados, confortables y seguros donde el usuario sienta absoluta reserva y confidencialidad respecto a sus padecimientos íntimos. \\
\addlinespace
\textbf{Corporalidad} \newline (\textit{Cuerpo vivido}) & 
Delicadeza y respeto al pudor corporal al desvestirse o ser examinado; tacto suave en inyecciones y curaciones de enfermería; respuesta afectuosa ante el llanto del niño o los dolores articulares del anciano. & 
El cuerpo doliente (\textit{nanay}) no es un objeto mecánico de intervención, sino el asiento directo de la vulnerabilidad ontológica. La calidad del cuidado reside en la delicadeza física, la calidez manual y la empatía somática que reconforta el cuerpo enfermo (\textit{onqoy}). \\
\bottomrule
\end{tabularx}
\end{table}

Esta articulación multidimensional confirma que, para el usuario del Centro de Salud Belén, la ``calidad de atención'' trasciende ampliamente los parámetros gerenciales o el abastecimiento instrumental; constituye una experiencia existencial totalizante donde la relación interpersonal, el tiempo dedicado, el espacio íntimo y el respeto al cuerpo configuran la esencia del cuidado digno y humanizado en el sistema público de salud.

\newpage
